\documentclass{article}
\usepackage[UTF8]{ctex}
\usepackage{amsmath,amssymb,amsfonts}
\usepackage{enumitem}
\usepackage{geometry}
\geometry{a4paper, margin=1in}
% === 错题本排版优化（不参与题目标识） ===
\setlength{\parindent}{0pt}
\setlength{\parskip}{0.35em plus 0.10em minus 0.08em}
\setlength{\emergencystretch}{3em}
\setlist{itemsep=0.25em, topsep=0.40em, parsep=0pt}
\allowdisplaybreaks[2]
% === 错题本排版优化结束 ===
\begin{document}
\section*{华中科技大学2024年高等代数真题}
\section*{123．华中科技大学2024年高等代数试卷}

\subsection{题目1}
1．计算行列式

\[
\left|\begin{array}{ccccc}
1 & 2 & 2 & \cdots & 2 \\
2 & 2 & 2 & \cdots & 2 \\
2 & 2 & 3 & \cdots & 2 \\
\vdots & \vdots & \vdots & & \vdots \\
2 & 2 & 2 & \cdots & 2024
\end{array}\right|
\]

\subsection{题目2}
2．给定线性方程组

\[
\left\{\begin{array}{l}
x_{1}+a x_{2}+a^{2} x_{3}=a^{3} \\
x_{1}+b x_{2}+b^{2} x_{3}=b^{3} \\
x_{1}+c x_{2}+c^{2} x_{3}=c^{3} \\
x_{1}+d x_{2}+d^{2} x_{3}=d^{3}
\end{array}\right.
\]

\begin{itemize}
  \item[（1）] 证明：当 \(a, b, c, d\) 互异时，方程组无解．
  \item[（2）] 当 \(a=b=c=-d\) 且 \(\xi=(1,1,-1)^{\prime}\) 为方程组的一个解时，求方程组的全部解。
\end{itemize}

\subsection{题目3}
3．设 \(A, B\) 均为 \(n\) 级实矩阵，且 \(A B=B A\) ，证明 \(\operatorname{det}\left(4 A^{2}+4 A B+5 B^{2}\right) \geq 0\) 。

\subsection{题目4}
4．设 \(A=\left(\begin{array}{ll}X & B \\ C & D\end{array}\right)\) ，其中 \(B, C, D\) 为给定的 \(n\) 级矩阵，\(X\) 是任意 \(n\) 级矩阵，当 \(X\) 变化时，求\(\operatorname{rank}(A)\) 的最小值．

\subsection{题目5}
5．设 \(V\) 是数域 \(F\) 上的2024维线性空间，\(\sigma_{1}, \sigma_{2}, \cdots, \sigma_{2024}\) 是 \(V\) 上两两不同的线性变换，证明存在 \(V\)的基 \(\alpha_{1}, \alpha_{2}, \cdots, \alpha_{2024}\) ，使得对任意的 \(1 \leq i \leq 2024\) ，都有 \(\sigma_{1}\left(\alpha_{i}\right), \sigma_{2}\left(\alpha_{i}\right), \cdots, \sigma_{2024}\left(\alpha_{i}\right)\) 互不相同．

\subsection{题目6}
6．已知 \(A\) 是 \(n\) 级实对称正定矩阵，\(B\) 是 \(n\) 级实对称矩阵。证明：\(B\) 是负定矩阵的充分必要条件是 \(A B\)的特征值均为负实数。

\subsection{题目7}
7．设 \(A\) 是3级实矩阵，满足对任意的 \(\alpha \in \mathbb{R}^{3}, \alpha\) 与 \(A \alpha\) 均正交、证明：存在向量 \(\beta \in \mathbb{R}^{3}\) ，使得\(A \alpha=\alpha \times \beta\) 。

\subsection{题目8}
8．已知 \(A, B\) 均为3级矩阵，证明： \(\operatorname{det}(A B-B A)=\frac{1}{3} \operatorname{tr}\left((A B-B A)^{3}\right)\) 。
\end{document}
